% Optional sharpened transfer fragment. It is distinct from frozen costs.
\section{Exact variation of the common-normal payoff difference}
\label{sec:raw-exact-variation}

\begin{proposition}[Exact variation]
\label{prop:raw-exact-variation}
For $\Delta\geq0$, $m\in\mathbb R$ and $\lambda>0$, let
$\ell=\log(101/100)$ and
$h(n)=f(m+\Delta+\lambda n)-f(m+\lambda n)$. Then
\[
 \operatorname{TV}(h)
 =2\min\{1,101(1-e^{-\Delta})\}.
\]
Consequently, the normal CDF transfer obeys the sufficient conditional
bound
\[
 |\mathbb E h(\widetilde N)-\mathbb E h(N)|
 \leq2D_N\min\{1,101(1-e^{-\Delta})\}.
\]
\end{proposition}
\begin{proof}
An increasing affine change of the independent variable preserves total
variation, so it suffices to study
$g(x)=f(x+\Delta)-f(x)$. If $\Delta=0$, the function is identically
zero. Suppose first that $0<\Delta<\ell$. Its explicit pieces are
\[
 g(x)=\begin{cases}
 0,&x\leq-\Delta,\\
 100(e^{x+\Delta}-1),&-\Delta<x\leq0,\\
 100e^x(e^\Delta-1),&0<x<\ell-\Delta,\\
 101-100e^x,&\ell-\Delta\leq x<\ell,\\
 0,&x\geq\ell.
 \end{cases}
\]
The pieces agree at every boundary. The first nonzero piece is increasing;
the middle piece is increasing; the last nonzero piece is decreasing.
The maximum at $x=\ell-\Delta$ is
$101-100e^{\ell-\Delta}=101(1-e^{-\Delta})<1$.
Because the continuous function increases from zero to this maximum and
then decreases to zero, its variation is twice its maximum.

If $\Delta\geq\ell$, the nonzero pieces instead are
\[
 g(x)=\begin{cases}
 100(e^{x+\Delta}-1),&-\Delta<x<\ell-\Delta,\\
 1,&\ell-\Delta\leq x\leq0,\\
 101-100e^x,&0<x<\ell.
 \end{cases}
\]
Outside their union the function is zero. The plateau has height one;
for $\Delta=\ell$ it reduces to a single point. The variation is two.
At $\Delta=\ell$, $101(1-e^{-\ell})=1$, so the two formulae agree.
The CDF estimate follows from the same signed Stieltjes integration by
parts used for the original variation upper bound, replacing its bound
two by the exact variation just derived.
\end{proof}

This sharpening is conditional on the endpoint gap. A uniform certificate
may replace $\Delta$ by a valid upper bound because
$\min\{1,101(1-e^{-\Delta})\}$ is nondecreasing. If that upper bound
exceeds $\ell$, the uniform estimate is still $2D_N$. For the
smoothing-first certificate, the gap in the CDF step is the represented
nonnegative $\widehat\Delta$, so a bound for the true gap alone does not
justify a tighter represented-gap transfer without its arithmetic
enlargement. This optional sharpening has not been used in the frozen
pointwise resource results. Any resource claim which uses it requires a
separately recorded financial certificate, prior and estimator rebuild.
